% fig_collision_box.tex — Fig. 4.2: a collision witness, schematic in a coordinate plane. Two
% baseline unit cells of the eighth grid overlap in an interval of length >= 1/8 in each
% coordinate; eroding both by rho = 1/100 leaves a core box of side >= 21/200 inside both solids.
% Used by sec5_registration.tex.
\begin{tikzpicture}[x=1cm,y=1cm,font=\small,>={Latex[length=1.6mm]}]
  % two unit cells on the eighth grid, offset by (1/8, 3/8) (drawn at 3.2 cm per unit)
  \def\S{3.2}
  \draw[gray!40,step=0.4] (-0.2,-0.2) grid (4.6,4.6);
  \draw[thick,blue!70!black] (0,0) rectangle (\S,\S);
  \draw[thick,red!70!black] (0.4,1.2) rectangle (0.4+\S,1.2+\S);
  \fill[black!15] (0.4+0.032,1.2+0.032) rectangle (\S-0.032,\S-0.032);
  \draw[very thick] (0.4+0.032,1.2+0.032) rectangle (\S-0.032,\S-0.032);
  \node[blue!70!black,anchor=south west] at (0.05,0.05) {cell of the other tile};
  \node[red!70!black,anchor=north east] at (0.4+\S-0.05,1.2+\S-0.05) {cell of the option $W$};
  \node[anchor=west,align=left] at (4.9,2.6) {overlap of the two unit cells:\\ each coordinate interval has length $\ge 1/8$\\ (all endpoints are eighth-integral)};
  \node[anchor=west,align=left] at (4.9,1.2) {shaded: both cells eroded by $\rho=1/100$,\\ side $\ge 1/8-2/100 = 21/200$; the eroded\\ cells lie inside the two solids, so the\\ solids overlap in a box of that side};
  \node[anchor=north,align=center] at (2.3,-0.4) {Schematic coordinate-plane section of the overlap estimate\\ (the certificate stores the actual boxes in integer units of $1/400$)};
\end{tikzpicture}
